\chapter{Requirement Analysis}
\label{chap:requirement-analysis}

\section{Stakeholder Analysis}
\label{section:stakeholder-analysis}
% <TIP: List your stakeholders for your project here/>

% Stakeholders are individuals, groups, or entities that
% have an interest, concern, or stake in a particular project, decision,
% organization, or system. These are individuals or groups who can affect or be
% affected by the outcomes of your project.
\subsection{Direct Users / End-Users}
\begin{itemize}
    \item \textbf{Law Students:} Need a realistic, interactive environment to study and practice civil procedure. They are directly affected by the UI complexity and accuracy of the courtroom flow.
    \item \textbf{Lawyers and Legal Practitioners:} Require an accurate tool to evaluate civil case winning potential based on facts. Their adoption depends on the simulator's alignment with the Civil and Commercial Code.
    \item \textbf{Judiciary Trainees and Court Staff:} Need to practice managing civil trials and drafting judgments. They require high-fidelity simulations of actual courtroom environments.
    \item \textbf{Litigants:} Want to simulate real-world disputes to understand their legal standing before hiring a lawyer. They represent a non-technical user base that relies on simplified legal terminology.
\end{itemize}

\subsection{External Influencers}
\begin{itemize}
    \item \textbf{Lawyers Council of Thailand:} Concerned with maintaining ethical standards and preventing the unauthorized practice of law. Their public stance could heavily influence professional adoption.
    \item \textbf{Office of the Judiciary / Court Administration:} Govern court procedures and are interested in modernizing legal training. The software's workflows must strictly align with their official administrative and procedural rules.
\end{itemize}

\subsection{Beneficiaries or Indirect Users}
\begin{itemize}
    \item \textbf{Law Firm Clients:} Desire successful and cost-effective resolution of their disputes. Their financial future can be influenced by the lawyer's use of the simulator.
    \item \textbf{The Justice System:} Aims to reduce court backlogs. Benefited indirectly if the tool helps users realize a case is weak, thereby discouraging frivolous lawsuits from being filed.
\end{itemize}


\section{User Stories and Use Cases}
\label{section:user-stories-use-cases}
% <TIP: For each stakeholder, write user stories that capture what
% they need from the system: "As a [stakeholder], I want to [action] so that [benefit]."
% For your most important user stories, expand them into use cases.
% A use case is a detailed description of how a system
% interacts with an external entity (such as a user or another system) to
% accomplish a specific goal. Use cases provide a high-level view of the
% functionality of a system and help in capturing and documenting its
% requirements from the perspective of end users.

% Make sure to use the
% appropriate type of use case for each scenario (brief, casual, and fully-dressed
% use case). />
\subsection{User Stories}

\subsubsection{Law Students}
\begin{itemize}
    \item As a law student, I want to manually control the Prosecutor Agent in a simulated trial so that I can practice my case-building, argumentation, and cross-examination dynamically.
    \item As a law student, I want to input evidence and conversation dynamically during a live mock trial so that I can learn how new information shifts the trial's direction in real-time.
    \item As a law student, I want to raise procedural objections against the MARL-trained Defender during the trial so that I can test my knowledge of Thai civil procedure rules.
\end{itemize}

\subsubsection{Lawyers and Legal Practitioners}
\begin{itemize}
    \item As a legal practitioner, I want to run a hybrid simulation against a MARL-trained Defender Agent so that I can stress-test my case arguments against realistic, auto-generated counter-arguments.
    \item As a legal practitioner, I want the Case Predictor Agent to analyze the completed trial simulation and output a win probability percentage so that I can evaluate case viability before accepting a client.
    \item As a legal practitioner, I want the system to automatically suggest relevant Supreme Court precedents during the post-trial feedback phase so that I can strengthen my real-world legal strategy.
\end{itemize}

\subsubsection{Judiciary Trainees and Court Staff}
\begin{itemize}
    \item As a judiciary trainee, I want to act as the manual Judge Agent while the AI handles the Prosecutor and Defender roles so that I can practice ruling on procedural objections and managing trial flow.
    \item As a judiciary trainee, I want to review the post-trial analysis generated by the Case Predictor so that I can compare my manual trial rulings against the system's objective legal evaluation.
\end{itemize}

\subsubsection{Litigants / General Public}
\begin{itemize}
    \item As a litigant, I want plain-language guidance prompts during simulation setup so that I can input my Thai Civil Law dispute details without deciphering complex legal jargon.
    \item As a litigant, I want the Case Predictor Agent to provide a simplified summary of potential legal outcomes post-simulation so that I can decide if it is worth hiring a lawyer.
    \item As a litigant, I want the system to suggest alternative dispute resolution (ADR) options based on the simulation results so that I can explore cheaper ways to settle my conflict outside of court.
\end{itemize}

\subsection{Use Cases}

This section describes the main use cases of the ChatreeDeka system, including the primary actors, objectives, preconditions, main flow, alternative flow, and postconditions.

% \textbf{Fully-Dressed Use Case: Evaluate Civil Case Winning Potential} \\
% \textbf{Primary Actor:} Lawyer / Legal Practitioner \\
% \textbf{Goal:} To determine the viability and likely outcome of a civil dispute based on the Civil and Commercial Code. \\
% \textbf{Preconditions:} The user is authenticated and has gathered basic facts and evidence from a prospective client. \\
% \textbf{Postconditions:} The system generates a detailed probability report with relevant legal precedents. \\
% \textbf{Main Success Scenario:}
% \begin{enumerate}
%     \item The user initiates a ``New Case Evaluation'' from the dashboard.
%     \item The system prompts the user to select the civil law category (e.g., Breach of Contract, Tort, Property).
%     \item The user inputs the specific facts, timeline of events, and available evidence using the structured input matrix.
%     \item The user submits the data for analysis.
%     \item The system's legal rule engine processes the inputs against the simulated Civil Code database.
%     \item The system calculates the strength of the case and identifies matching historical precedents.
%     \item The system displays a comprehensive evaluation report, highlighting strengths, weaknesses, and a percentage-based probability of success.
% \end{enumerate}
% \textbf{Extensions:}
% \begin{itemize}
%     \item \textit{3a. Incomplete Data:} If the user fails to input mandatory elements, the system flags the missing fields and pauses the evaluation.
%     \item \textit{5a. No Precedent Match:} If the system cannot find a matching precedent, it outputs an evaluation based strictly on statutory interpretation.
% \end{itemize}

% \textbf{Casual Use Case: Match Supreme Court Precedents} \\
% \textbf{Primary Actor:} Lawyer / Legal Practitioner \\
% \textbf{Scenario:} While reviewing a simulated case evaluation, the lawyer clicks the ``Find Precedents'' button. The system extracts key legal tags (e.g., ``vicarious liability,'' ``agency'') from the user's input and queries the internal database of historical rulings. The system returns a curated list of case summaries that match the current scenario, allowing the lawyer to add these citations directly to their simulated case file.

% \subsubsection{Law Students}

% \textbf{Fully-Dressed Use Case: Submit Legal Pleading for Validation} \\
% \textbf{Primary Actor:} Law Student \\
% \textbf{Goal:} To draft a correct initial complaint or statement of defense and receive feedback. \\
% \textbf{Preconditions:} The student has selected a mock case and is assigned a role (Plaintiff/Defendant). \\
% \textbf{Postconditions:} The pleading is accepted into the simulation and feedback is logged. \\
% \textbf{Main Success Scenario:}
% \begin{enumerate}
%     \item The student navigates to the ``Document Drafting'' module.
%     \item The system provides a blank template appropriate for the required pleading.
%     \item The student types their legal arguments, citing specific sections of the Civil and Commercial Code.
%     \item The student submits the document to the virtual court.
%     \item The system scans the document for required structural elements (e.g., jurisdiction, relief sought) and validates the cited laws against the case facts.
%     \item The system accepts the document and unlocks the next phase of the trial.
%     \item The system provides a sidebar with automated feedback on the logical structure and statutory accuracy of the pleading.
% \end{enumerate}
% \textbf{Extensions:}
% \begin{itemize}
%     \item \textit{5a. Fatal Flaw Detected:} If the student cites a completely irrelevant law or misses a mandatory procedural element, the system rejects the filing, highlights the error, and prompts a rewrite.
% \end{itemize}

% \textbf{Brief Use Case: Raise Procedural Objection} \\
% \textbf{Primary Actor:} Law Student \\
% \textbf{Scenario:} During the virtual witness examination phase, the student clicks the ``Objection'' hotkey. The simulation pauses, and the student selects the grounds for objection (e.g., ``Hearsay,'' ``Leading Question'') from a dropdown menu. The virtual judge immediately sustains or overrules the objection based on the predefined rules of the scenario, adjusting the trial score accordingly.

% \subsubsection{Judiciary Trainees and Court Staff}

% \textbf{Casual Use Case: Draft and Issue a Civil Judgment} \\
% \textbf{Primary Actor:} Judiciary Trainee \\
% \textbf{Scenario:} After all evidence and arguments have been submitted in a simulated trial, the trainee opens the ``Judgment Drafting'' interface. The system displays a split screen: one side contains the full record of the simulated trial, and the other contains a text editor. As the trainee writes the ruling, the system's auto-suggest feature recommends standard legal phrasing and statutory citations. Once finalized, the trainee clicks ``Issue Judgment,'' which closes the simulation and generates a final PDF for grading.

% \textbf{Brief Use Case: Rule on Evidence Admissibility} \\
% \textbf{Primary Actor:} Judiciary Trainee \\
% \textbf{Scenario:} A simulated plaintiff attempts to submit a document into evidence. The system prompts the trainee on the judge's dashboard to review the document. The trainee clicks ``Admit'' or ``Exclude'' based on their assessment of its relevance and authenticity. The system logs the decision and proceeds with the trial.

% \subsubsection{Litigants / General Public}

% \textbf{Casual Use Case: Complete Plain-Language Legal Assessment} \\
% \textbf{Primary Actor:} Litigant / General Public \\
% \textbf{Scenario:} A user accesses the public portal and selects ``Evaluate My Problem.'' The system presents a conversational, wizard-style questionnaire. Instead of asking for legal elements, it asks factual questions like ``Did you sign a piece of paper?'' or ``Were you physically injured?'' Once the user finishes the flow, the system translates these answers into a basic legal assessment, explaining their situation in simple terms without using statutory jargon.

% \textbf{Brief Use Case: Generate Mediation Recommendation} \\
% \textbf{Primary Actor:} Litigant / General Public \\
% \textbf{Scenario:} After completing a legal assessment, the system calculates that the user's financial claim is very small or the case is highly risky. The system automatically displays a ``Mediation Recommended'' alert, providing a brief explanation of how out-of-court settlement works and offering a link to local dispute resolution resources.
\subsubsection{UC-01: Input Case Details and Evidence}

\begin{longtable}{|p{4cm}|p{10cm}|}
\hline
\textbf{Use Case ID} & UC-01 \\ \hline
\textbf{Use Case Name} & Input Case Details and Evidence \\ \hline
\textbf{Primary Actor} & Law Student, Lawyer / Legal Practitioner, Judiciary Trainee / Court Staff, Litigant / General Public \\ \hline
\textbf{Description} & The user enters case facts and supporting evidence into the system before starting a simulation. \\ \hline
\textbf{Preconditions} & The user has opened the case setup page. \\ \hline
\textbf{Trigger} & The user selects Start New Case. \\ \hline
\textbf{Main Flow} & 
1. The user opens the case setup page. \newline
2. The user enters dispute details. \newline
3. The user uploads or types evidence and facts. \newline
4. The system validates the input. \newline
5. The system stores the case information for simulation. \\ \hline
\textbf{Alternative Flow} & 
If required information is missing, the system asks the user to complete it before proceeding. \\ \hline
\textbf{Postconditions} & 
A case draft is created and ready for role configuration. \\ \hline
\end{longtable}

\subsubsection{UC-02: Select Role Configuration}

\begin{longtable}{|p{4cm}|p{10cm}|}
\hline
\textbf{Use Case ID} & UC-02 \\ \hline
\textbf{Use Case Name} & Select Role Configuration \\ \hline
\textbf{Primary Actor} & Law Student, Lawyer / Legal Practitioner, Judiciary Trainee / Court Staff \\ \hline
\textbf{Description} & The user selects which courtroom roles are controlled manually and which are controlled by AI before starting the simulation. \\ \hline
\textbf{Preconditions} & Case details and evidence have already been entered into the system. \\ \hline
\textbf{Trigger} & The user proceeds from case setup to simulation settings. \\ \hline
\textbf{Main Flow} & 
1. The user opens the role configuration page. \newline
2. The system displays the available roles, such as judge, prosecutor, and defender. \newline
3. The user chooses whether each role is controlled by a human or AI. \newline
4. The system checks whether the selected configuration is valid. \newline
5. The system saves the simulation settings. \\ \hline
\textbf{Alternative Flow} & 
If the selected role combination is incomplete or invalid, the system asks the user to revise the configuration. \\ \hline
\textbf{Postconditions} & 
The role configuration is saved and the system is ready to start the trial simulation. \\ \hline
\end{longtable}

\subsubsection{UC-03: Start Trial Simulation}

\begin{longtable}{|p{4cm}|p{10cm}|}
\hline
\textbf{Use Case ID} & UC-03 \\ \hline
\textbf{Use Case Name} & Start Trial Simulation \\ \hline
\textbf{Primary Actor} & Law Student, Lawyer / Legal Practitioner, Judiciary Trainee / Court Staff, Litigant / General Public \\ \hline
\textbf{Description} & The user starts the trial simulation using the entered case details and selected role configuration. \\ \hline
\textbf{Preconditions} & Case details have been entered and role configuration has been completed. \\ \hline
\textbf{Trigger} & The user clicks the Start Simulation button. \\ \hline
\textbf{Main Flow} & 
1. The user selects Start Simulation. \newline
2. The system loads the case data and role configuration. \newline
3. The system initializes the trial environment and participating roles. \newline
4. The simulation begins with the first courtroom interaction. \newline
5. The system records all interactions in the trial transcript. \\ \hline
\textbf{Alternative Flow} & 
If required case data or role settings are missing, the system prevents the simulation from starting and asks the user to complete the missing information. \\ \hline
\textbf{Postconditions} & 
The trial simulation session is active and ready for user interaction. \\ \hline
\end{longtable}

\subsubsection{UC-04: Inject Evidence and Add New Facts}

\begin{longtable}{|p{4cm}|p{10cm}|}
\hline
\textbf{Use Case ID} & UC-04 \\ \hline
\textbf{Use Case Name} & Inject Evidence and Add New Facts \\ \hline
\textbf{Primary Actor} & Law Student, Lawyer / Legal Practitioner, Litigant / General Public \\ \hline
\textbf{Description} & The user introduces additional evidence or new facts into the case while the trial simulation is in progress. \\ \hline
\textbf{Preconditions} & A trial simulation is currently running. \\ \hline
\textbf{Trigger} & The user selects the option to add evidence or new facts during the simulation. \\ \hline
\textbf{Main Flow} & 
1. The user opens the evidence injection interface. \newline
2. The user enters new facts or uploads supporting evidence. \newline
3. The system validates and records the newly submitted information. \newline
4. The system updates the simulation context with the new information. \newline
5. The AI or manual participants continue the simulation based on the updated case state. \\ \hline
\textbf{Alternative Flow} & 
If the submitted information is incomplete, irrelevant, or unreadable, the system asks the user to correct or clarify it before acceptance. \\ \hline
\textbf{Postconditions} & 
The case state is updated with the new evidence or facts, and the simulation continues with the revised information. \\ \hline
\end{longtable}

\subsubsection{UC-05: Raise Procedural Objection}

\begin{longtable}{|p{4cm}|p{10cm}|}
\hline
\textbf{Use Case ID} & UC-05 \\ \hline
\textbf{Use Case Name} & Raise Procedural Objection \\ \hline
\textbf{Primary Actor} & Law Student, Lawyer / Legal Practitioner \\ \hline
\textbf{Description} & The user raises a procedural objection during the trial simulation when they detect an improper statement, action, or evidence submission. \\ \hline
\textbf{Preconditions} & A trial simulation is active and the current stage allows objections. \\ \hline
\textbf{Trigger} & The user selects the objection option during an active exchange in the simulation. \\ \hline
\textbf{Main Flow} & 
1. The user selects Raise Objection. \newline
2. The system displays the available objection types. \newline
3. The user chooses the objection type and submits it. \newline
4. The system records the objection in the trial transcript. \newline
5. The objection is forwarded to the judge role for consideration or ruling. \\ \hline
\textbf{Alternative Flow} & 
If objections are not allowed at the current point in the simulation, the system rejects the request and informs the user. \\ \hline
\textbf{Postconditions} & 
The procedural objection is recorded and awaits a ruling or resolution within the simulation. \\ \hline
\end{longtable}

\subsubsection{UC-06: View Post-Trial Analysis and Feedback}

\begin{longtable}{|p{4cm}|p{10cm}|}
\hline
\textbf{Use Case ID} & UC-06 \\ \hline
\textbf{Use Case Name} & View Post-Trial Analysis and Feedback \\ \hline
\textbf{Primary Actor} & Law Student, Lawyer / Legal Practitioner, Judiciary Trainee / Court Staff, Litigant / General Public \\ \hline
\textbf{Description} & The user reviews the analysis generated by the system after the trial simulation, including feedback on legal reasoning, case strength, and possible outcomes. \\ \hline
\textbf{Preconditions} & The trial simulation has been completed. \\ \hline
\textbf{Trigger} & The user opens the post-trial results page after finishing the simulation. \\ \hline
\textbf{Main Flow} & 
1. The user opens the post-trial analysis page. \newline
2. The system analyzes the trial transcript, submitted evidence, and simulation outcome. \newline
3. The system generates feedback and case analysis results. \newline
4. The system displays the win probability and relevant feedback to the user. \newline
5. The user reviews the analysis for learning, strategy, or decision-making purposes. \\ \hline
\textbf{Alternative Flow} & 
If the analysis cannot be fully generated due to incomplete trial data, the system displays partial feedback and informs the user of the limitation. \\ \hline
\textbf{Postconditions} & 
The user receives post-trial insights and can use them for further review, preparation, or decision-making. \\ \hline
\end{longtable}

\section{Use Case Diagrams}
\label{section:use-case-diagrams}
\begin{figure}[H]
    \centering
    \includegraphics[width=0.9\textwidth]{images/chartee_usecase.png}
    \caption{Use Case Diagram for ChatreeDeka}
\end{figure}

\section{User Interface Design}
\label{section:user-interface-design}
% <TIP: Put the initial design of your application here. You can
% showcase a detailed design of a specific page or a sitemap of your application.
% See an example below./>

\begin{figure}[h]
    \centering
    \includegraphics[width=0.8\textwidth]{images/ui-presim.png}
    \caption{User Interface Design - Pre-Simulation Page}
\end{figure}

\begin{figure}[h]
    \centering
    \includegraphics[width=0.8\textwidth]{images/ui-sim.png}
    \caption{User Interface Design - Simulation Page}
\end{figure}

\begin{figure}[h]
    \centering
    \includegraphics[width=0.8\textwidth]{images/ui-postsim.png}
    \caption{User Interface Design - Post-Simulation Page}
\end{figure}
